\icode{E-MD-SI-L2-260625-001}
2023年诺贝尔物理学奖授予了“为研究物质中的电子动力学而产生阿秒（\qty{e-18}{s}）光脉冲实验方法”的三位物理学家，他们的实验为人类探索原子和分子内部的电子世界提供了新的工具。科学研究表明，人类所能测量的最短长度和最短时间都存在一个极限，分别称为普朗克长度\(l_p\)和普朗克时间\(t_p\)，两者之间满足关系\(l_p = t_p\cdot c\)（\(c\)为真空中的光速）。推导普朗克长度和普朗克时间需要用到三个常数：万有引力常量\(G\)、光速\(c\)和普朗克常数\(h\)。已知\(h\)的单位为\qty{}{J\cdot s}，下列关系式正确的是\paren[]
\begin{choices}
    \item \(l_p = \frac{Gh}{c}\)
    \item \(l_p = \sqrt{\frac{Gh}{c}}\)
    \item \(l_p = \sqrt{\frac{Gh}{c^3}}\)
    \item \(l_p = \sqrt{\frac{Gh}{c^5}}\)
\end{choices}
\begin{answer}
    C\qquad 提示：
\end{answer}